\section{Artifact contract and ObserverCards}

The workbench needs a reusable boundary if researchers are to evaluate their
own observers without rebuilding each fixture. ObserverBench v0 exposes a small
Python composition contract with version identifier
\texttt{observerbench.api.v0}. It separates four responsibilities:

\begin{itemize}
  \item a state estimator returns the task-defined estimate;
  \item an actuation-direction provider returns a direction but not an action
  magnitude;
  \item a controller maps target and estimate to a scalar action;
  \item a task owns the plant, intervention rule, controller settings, and
  metrics.
\end{itemize}

The task registry reports which bundled tasks accept this interface. At v0,
only analytic Ctl-1 has a tested external-observer adapter. Its benchmark-facing
wrapper constructs the controller internally and reuses the existing data
generator, readouts, normalization, and metrics. A lower-level composition may
use a custom controller, but the output is marked non-comparable. Runs use
strict JSON and record the contract and result-schema versions, component
classes and parameters, package and dependency versions, source revision,
configuration and result hashes, and caller-supplied observer provenance.

This is a bounded claim. The contract proves that an outside observer can run
unchanged against one existing fixture. The trained and IOI tasks remain fixed
reproduction tasks. v0 does not provide dynamic plugin discovery, a fitting or
split lifecycle, remote execution, or a universal bring-your-own-model API.

ObserverCards summarize the scientific result. A card records the task,
observer components, access regime, estimand, measurement design, validation
target, main metrics, detected failure modes, recommendation, and scope limits.
Recommendations are metric-derived rather than observer-name-derived. Two
observer names with the same metrics receive the same recommendation; the same
name can receive different recommendations under different intervention
distributions.

For Ctl-2, the card now distinguishes estimator from direction, reports the
observer self-gain and conditioning diagnostics, and uses ``sign-incompatible
feedback'' or ``sustained target-error growth'' rather than unqualified
``instability.'' For the IOI random-subset task, it can recommend the additive
observer. For the primary-stratified task, its final recommendation waits for
the capacity-matched pair analysis.
